\section{Validação de Consistência dos Dados}

\subsection{Quantidade de Registros}

A tabela abaixo resume a quantidade de registros encontrados nos logs de telemetria antes e depois do processo de deduplicação (remoção de duplicados exatos).

\begin{table}[H]
\IBGEtab{%
  \caption{Quantidade de Registros}%
  \label{tab:data_validacao_de_dados_1}%
}{%
  \resizebox{\textwidth}{!}{%
    \begin{tabular}{lcccc}
      \toprule
      Fabricante & Registros Originais & Registros Duplicados Removidos & Registros Limpos & \% de Redundância \\
      \midrule
      \textbf{Ruckus} & 318.074 & 0 & 318.074 & 0,00\% \\
      \textbf{UniFi} & 36.180 & 0 & 36.180 & 0,00\% \\
      \bottomrule
    \end{tabular}%
  }%
}{%
  \fonte{Produzido pelos autores.}%
}
\end{table}

\textit{Nota: Registros duplicados ocorrem por sobreposição em polling repetido do coletor. A remoção evita viés nas análises estatísticas.}

\subsection{Período Coberto pela Coleta}

\begin{table}[H]
\IBGEtab{%
  \caption{Período Coberto pela Coleta}%
  \label{tab:data_validacao_de_dados_2}%
}{%
  \resizebox{\textwidth}{!}{%
    \begin{tabular}{lccc}
      \toprule
      Fabricante & Data/Hora Inicial (Min) & Data/Hora Final (Max) & Duração Efetiva \\
      \midrule
      \textbf{Ruckus} & 23/06/2026 00:00 & 10/07/2026 23:55 & 17 days 23:55:00 \\
      \textbf{UniFi} & 23/06/2026 00:02 & 10/07/2026 23:57 & 17 days 23:55:00 \\
      \bottomrule
    \end{tabular}%
  }%
}{%
  \fonte{Produzido pelos autores.}%
}
\end{table}

\textit{Nota: Ambos os datasets cobrem o mesmo período temporal (de 23/06/2026 00:00 a 10/07/2026 23:55), garantindo que os cenários de carga de rede sejam diretamente comparáveis.}

\subsection{Dados Ausentes (Missing Values)}

O percentual de dados nulos ou ausentes para cada variável após a limpeza:

\begin{table}[H]
\IBGEtab{%
  \caption{Dados Ausentes (Missing Values)}%
  \label{tab:data_validacao_de_dados_3}%
}{%
  \resizebox{\textwidth}{!}{%
    \begin{tabular}{lccc}
      \toprule
      Coluna Padronizada & \% Nulos Ruckus & \% Nulos UniFi & Tipo de Dado \\
      \midrule
      \texttt{timestamp} & 0,00\% & 0,00\% & Datetime \\
      \texttt{AP} & 0,00\% & 0,00\% & Categórico \\
      \texttt{cliente (MAC)} & 0,00\% & 0,00\% & Categórico \\
      \texttt{RSSI} & 10,34\% & 0,00\% & Numérico \\
      \texttt{signal} & 10,34\% & 0,00\% & Numérico \\
      \texttt{noise} & 10,34\% & 0,00\% & Numérico \\
      \texttt{canal} & 0,00\% & 0,00\% & Categórico \\
      \texttt{throughput\_tx\_mbps} & 0,00\% & 0,00\% & Numérico \\
      \texttt{throughput\_rx\_mbps} & 0,00\% & 0,00\% & Numérico \\
      \texttt{link\_speed\_tx\_mbps} & 6,51\% & 0,00\% & Numérico \\
      \texttt{link\_speed\_rx\_mbps} & 100,00\% & 0,00\% & Numérico \\
      \texttt{retransmissoes\_pct} & 0,00\% & 0,00\% & Numérico \\
      \texttt{volume\_tx\_mb} & 0,00\% & 0,00\% & Numérico \\
      \texttt{volume\_rx\_mb} & 0,00\% & 0,00\% & Numérico \\
      \texttt{padrão} & 0,00\% & 0,00\% & Categórico \\
      \bottomrule
    \end{tabular}%
  }%
}{%
  \fonte{Produzido pelos autores.}%
}
\end{table}

\subsection{Tratamento de Valores Nulos}

\begin{itemize}
  \item \textbf{Link Speed no Ruckus}: A taxa de modulação física de TX (\texttt{link\_speed\_tx\_mbps}) foi extraída com sucesso a partir do campo \texttt{throughput\_rate} (em Kbps) nos logs de telemetria da controladora SmartZone. A taxa RX correspondente (\texttt{link\_speed\_rx\_mbps}) não é informada nos logs brutos e é mantida como \texttt{NaN}.
  \item \textbf{Ruído (Noise) e Sinal (Signal) em Ruckus}: Alguns registros continham valores de sinal zerados ou vazios. Onde o sinal era -100 (vazio padrão) e SNR era 0, o ruído foi definido como NaN para evitar distorções de cálculo.
  \item \textbf{Valores Nulos na UniFi}: O dataset UniFi apresenta consistência total (0\% de nulos nas métricas principais de telemetria).

\end{itemize}
\subsection{Padronização das Unidades de Medida}

As colunas foram padronizadas da seguinte forma:
\begin{enumerate}
  \item \textbf{Sinal e Ruído}: Potência medida em \textbf{dBm}.
  \item \textbf{RSSI/SNR}: Medido em \textbf{dB} (ou índice relativo na UniFi).
  \item \textbf{Throughput}: Convertido em \textbf{Mbps} (Megabits por segundo) para Downstream (\texttt{tx}) e Upstream (\texttt{rx}).
  \item \textbf{Volume de Dados}: Convertido em \textbf{MB} (Megabytes) acumulados por cliente no período de observação.
  \item \textbf{Retransmissões}: Medidas em porcentagem (\textbf{\%}) de pacotes retransmitidos.
  \item \textbf{Link Speed}: Modulação de velocidade física convertida em \textbf{Mbps}.

\end{enumerate}
\subsection{Equivalência de Métricas (Mapeamento Ruckus vs UniFi)}

A tabela abaixo mostra a relação entre os campos extraídos de cada fabricante que representam as mesmas grandezas físicas:

\begin{table}[H]
\IBGEtab{%
  \caption{Equivalência de Métricas (Mapeamento Ruckus vs UniFi)}%
  \label{tab:data_validacao_de_dados_4}%
}{%
  \resizebox{\textwidth}{!}{%
    \begin{tabular}{lccc}
      \toprule
      Grandeza Física & Métrica Ruckus (Raw) & Métrica UniFi (Raw) & Coluna Padronizada Final \\
      \midrule
      Potência do Sinal & \texttt{signal\_rssi} & \texttt{signal} & \texttt{signal} (dBm) \\
      Relação Sinal/Ruído & \texttt{snr} & \texttt{rssi} & \texttt{RSSI} (dB) \\
      Throughput Downstream & \texttt{tx\_bytes / uptime} & \texttt{tx\_bytes\_r \textsuperscript{*} 8} & \texttt{throughput\_tx\_mbps} \\
      Throughput Upstream & \texttt{rx\_bytes / uptime} & \texttt{rx\_bytes\_r \textsuperscript{*} 8} & \texttt{throughput\_rx\_mbps} \\
      Retransmissões & \texttt{tx\_retries / (tx\_pkts + tx\_retries)} & \texttt{100 - (ccq / 10)} & \texttt{retransmissoes\_pct} \\
      Volume Trafegado & \texttt{tx\_bytes}, \texttt{rx\_bytes} & \texttt{tx\_bytes\_r \textit{ uptime}, \texttt{rx\_bytes\_r } uptime} & \texttt{volume\_tx\_mb}, \texttt{volume\_rx\_mb} \\
      Modulação de Velocidade & \texttt{throughput\_rate} & \texttt{tx\_rate}, \texttt{rx\_rate} & \texttt{link\_speed\_tx\_mbps}, \texttt{link\_speed\_rx\_mbps} \\
      \bottomrule
    \end{tabular}%
  }%
}{%
  \fonte{Produzido pelos autores.}%
}
\end{table}
